
\documentclass{article}
\usepackage{colortbl}
\usepackage{makecell}
\usepackage{multirow}
\usepackage{supertabular}

\begin{document}

\newcounter{utterance}

\centering \large Interaction Transcript for game `dond', experiment `coop\_en', episode 5 with Qwen3{-}32B{-}t0.0.
\vspace{24pt}

{ \footnotesize  \setcounter{utterance}{1}
\setlength{\tabcolsep}{0pt}
\begin{supertabular}{c@{$\;$}|p{.15\linewidth}@{}p{.15\linewidth}p{.15\linewidth}p{.15\linewidth}p{.15\linewidth}p{.15\linewidth}}
    \# & $\;$A & \multicolumn{4}{c}{Game Master} & $\;\:$B\\
    \hline

    \theutterance \stepcounter{utterance}  
    & & \multicolumn{4}{p{0.6\linewidth}}{
        \cellcolor[rgb]{0.9,0.9,0.9}{
            \makecell[{{p{\linewidth}}}]{
                \texttt{\tiny{[P1$\langle$GM]}}
                \texttt{You are playing a cooperative negotiation game in which you have to agree on how to divide a set of items among you and another player.} \\
\\ 
\texttt{Rules:} \\
\texttt{(a) You and the other player are given a set of items. Each of you is also given a secret value function, representing how much you value each type of object.} \\
\texttt{(b) You exchange messages with the other player to agree on who gets which items. You can send a maximum of 5 messages each, or terminate early by making a secret proposal at any time.} \\
\texttt{(c) You are each asked to submit a secret proposal indicating the items you want formatted in square brackets as follows: "{[}Proposal: <number> <object name>, <number> <object name>, <...>{]}"} \\
\texttt{(d) If your proposals are complementary, i.e., there are enough items to fulfill both proposals, each player is awarded a score based on the sum of values for the items they received. Otherwise, both of you get zero points.} \\
\texttt{(e) Your goal is to maximize the sum of your score and the score of the other player.} \\
\\ 
\texttt{Let us start.} \\
\\ 
\texttt{The set of available items is:} \\
\\ 
\texttt{2 chalks, 2 shoes, 2 fruits} \\
\\ 
\texttt{Your secret value function is:} \\
\\ 
\texttt{chalk: 4, shoe: 1, fruit: 0} \\
\\ 
\texttt{IMPORTANT: Your messages, unless it is the secret proposal, are directly transmitted to the other player, so do not include any response to the rules or text announcing your message. To make a secret proposal, use the indicated format. Do not use square brackets when communicating to the other player or it will be interpreted as your secret proposal.} \\
            }
        }
    }
    & \\ \\

    \theutterance \stepcounter{utterance}  
    & \multicolumn{4}{p{0.6\linewidth}}{
        \cellcolor[rgb]{0.9,0.9,0.9}{
            \makecell[{{p{\linewidth}}}]{
                \texttt{\tiny{[P1$\rangle$GM]}}
                \texttt{<think>} \\
\texttt{Okay, let's see. I need to divide 2 chalks, 2 shoes, and 2 fruits between me and another player. My value function is chalk:4, shoe:1, fruit:0. The other player's values are secret, so I don't know what they value. The goal is to maximize the total score (mine plus theirs) by making complementary proposals.} \\
\\ 
\texttt{First, I should figure out what the other player might value. Since fruits are worth 0 to me, maybe they value them. Chalks are high for me, so maybe they are low for them, or vice versa. Shoes are low for me, maybe they are higher for them.} \\
\\ 
\texttt{I need to communicate with the other player to find out their preferences. But I can only send 5 messages each. Let me start by asking them what they value. But wait, the messages are directly transmitted, so I can't include any meta{-}text. So I need to phrase it as a message they can understand.} \\
\\ 
\texttt{First message: "What do you value most? Let me know so we can split fairly." That's straightforward. Then, depending on their response, I can adjust my proposal.} \\
\\ 
\texttt{Suppose they reply that they value fruits the most. Then I can take the chalks and shoes, and they take the fruits. Since I don't value fruits, that's okay. But I need to check if their proposal is complementary. If they take all fruits, and I take the rest, that works. But maybe they also want some chalks or shoes.} \\
\\ 
\texttt{Alternatively, if they value shoes, maybe they want shoes and fruits. Then I can take the chalks and maybe one shoe. But I need to make sure the total items add up. Let's think about possible combinations.} \\
\\ 
\texttt{The total items are 2 of each. If I take 2 chalks and 2 shoes, that's 4 items, but there are only 6 total. Wait, no, there are 2 chalks, 2 shoes, 2 fruits. So total 6 items. Each player can take some combination. For example, if I take 2 chalks and 2 shoes, that's 4 items, leaving 2 fruits for the other. But if the other takes 2 fruits, that's 2 items. But maybe they want more. Alternatively, if I take 2 chalks and 1 shoe, leaving 1 shoe and} \\
            }
        }
    }
    & & \\ \\

    \theutterance \stepcounter{utterance}  
    & & & \multicolumn{2}{p{0.3\linewidth}}{
        \cellcolor[rgb]{0.9,0.9,0.9}{
            \makecell[{{p{\linewidth}}}]{
                \texttt{\tiny{[GM$|$GM]}}
                \texttt{continue} \\
            }
        }
    }
    & & \\ \\

    \theutterance \stepcounter{utterance}  
    & & & \multicolumn{2}{p{0.3\linewidth}}{
        \cellcolor[rgb]{0.9,0.9,0.9}{
            \makecell[{{p{\linewidth}}}]{
                \texttt{\tiny{[GM$|$GM]}}
                \texttt{<think>} \\
\texttt{Okay, let's see. I need to divide 2 chalks, 2 shoes, and 2 fruits between me and another player. My value function is chalk:4, shoe:1, fruit:0. The other player's values are secret, so I don't know what they value. The goal is to maximize the total score (mine plus theirs) by making complementary proposals.} \\
\\ 
\texttt{First, I should figure out what the other player might value. Since fruits are worth 0 to me, maybe they value them. Chalks are high for me, so maybe they are low for them, or vice versa. Shoes are low for me, maybe they are higher for them.} \\
\\ 
\texttt{I need to communicate with the other player to find out their preferences. But I can only send 5 messages each. Let me start by asking them what they value. But wait, the messages are directly transmitted, so I can't include any meta{-}text. So I need to phrase it as a message they can understand.} \\
\\ 
\texttt{First message: "What do you value most? Let me know so we can split fairly." That's straightforward. Then, depending on their response, I can adjust my proposal.} \\
\\ 
\texttt{Suppose they reply that they value fruits the most. Then I can take the chalks and shoes, and they take the fruits. Since I don't value fruits, that's okay. But I need to check if their proposal is complementary. If they take all fruits, and I take the rest, that works. But maybe they also want some chalks or shoes.} \\
\\ 
\texttt{Alternatively, if they value shoes, maybe they want shoes and fruits. Then I can take the chalks and maybe one shoe. But I need to make sure the total items add up. Let's think about possible combinations.} \\
\\ 
\texttt{The total items are 2 of each. If I take 2 chalks and 2 shoes, that's 4 items, but there are only 6 total. Wait, no, there are 2 chalks, 2 shoes, 2 fruits. So total 6 items. Each player can take some combination. For example, if I take 2 chalks and 2 shoes, that's 4 items, leaving 2 fruits for the other. But if the other takes 2 fruits, that's 2 items. But maybe they want more. Alternatively, if I take 2 chalks and 1 shoe, leaving 1 shoe and} \\
            }
        }
    }
    & & \\ \\

    \theutterance \stepcounter{utterance}  
    & & \multicolumn{4}{p{0.6\linewidth}}{
        \cellcolor[rgb]{0.9,0.9,0.9}{
            \makecell[{{p{\linewidth}}}]{
                \texttt{\tiny{[GM$\rangle$P2]}}
                \texttt{You are playing a cooperative negotiation game in which you have to agree on how to divide a set of items among you and another player.} \\
\\ 
\texttt{Rules:} \\
\texttt{(a) You and the other player are given a set of items. Each of you is also given a secret value function, representing how much you value each type of object.} \\
\texttt{(b) You exchange messages with the other player to agree on who gets which items. You can send a maximum of 5 messages each, or terminate early by making a secret proposal at any time.} \\
\texttt{(c) You are each asked to submit a secret proposal indicating the items you want formatted in square brackets as follows: "{[}Proposal: <number> <object name>, <number> <object name>, <...>{]}"} \\
\texttt{(d) If your proposals are complementary, i.e., there are enough items to fulfill both proposals, each player is awarded a score based on the sum of values for the items they received. Otherwise, both of you get zero points.} \\
\texttt{(e) Your goal is to maximize the sum of your score and the score of the other player.} \\
\\ 
\texttt{Let us start.} \\
\\ 
\texttt{The set of available items is:} \\
\\ 
\texttt{2 chalks, 2 shoes, 2 fruits} \\
\\ 
\texttt{Your secret value function is:} \\
\\ 
\texttt{chalk: 2, shoe: 0, fruit: 3} \\
\\ 
\texttt{IMPORTANT: Your messages, unless it is the secret proposal, are directly transmitted to the other player, so do not include any response to the rules or text announcing your message. To make a secret proposal, use the indicated format. Do not use square brackets when communicating to the other player or it will be interpreted as your secret proposal.} \\
\\ 
\\ 
\texttt{<think>} \\
\texttt{Okay, let's see. I need to divide 2 chalks, 2 shoes, and 2 fruits between me and another player. My value function is chalk:4, shoe:1, fruit:0. The other player's values are secret, so I don't know what they value. The goal is to maximize the total score (mine plus theirs) by making complementary proposals.} \\
\\ 
\texttt{First, I should figure out what the other player might value. Since fruits are worth 0 to me, maybe they value them. Chalks are high for me, so maybe they are low for them, or vice versa. Shoes are low for me, maybe they are higher for them.} \\
\\ 
\texttt{I need to communicate with the other player to find out their preferences. But I can only send 5 messages each. Let me start by asking them what they value. But wait, the messages are directly transmitted, so I can't include any meta{-}text. So I need to phrase it as a message they can understand.} \\
\\ 
\texttt{First message: "What do you value most? Let me know so we can split fairly." That's straightforward. Then, depending on their response, I can adjust my proposal.} \\
\\ 
\texttt{Suppose they reply that they value fruits the most. Then I can take the chalks and shoes, and they take the fruits. Since I don't value fruits, that's okay. But I need to check if their proposal is complementary. If they take all fruits, and I take the rest, that works. But maybe they also want some chalks or shoes.} \\
\\ 
\texttt{Alternatively, if they value shoes, maybe they want shoes and fruits. Then I can take the chalks and maybe one shoe. But I need to make sure the total items add up. Let's think about possible combinations.} \\
\\ 
\texttt{The total items are 2 of each. If I take 2 chalks and 2 shoes, that's 4 items, but there are only 6 total. Wait, no, there are 2 chalks, 2 shoes, 2 fruits. So total 6 items. Each player can take some combination. For example, if I take 2 chalks and 2 shoes, that's 4 items, leaving 2 fruits for the other. But if the other takes 2 fruits, that's 2 items. But maybe they want more. Alternatively, if I take 2 chalks and 1 shoe, leaving 1 shoe and} \\
            }
        }
    }
    & \\ \\

    \theutterance \stepcounter{utterance}  
    & & & \multicolumn{4}{p{0.6\linewidth}}{
        \cellcolor[rgb]{0.9,0.9,0.9}{
            \makecell[{{p{\linewidth}}}]{
                \texttt{\tiny{[GM$\langle$P2]}}
                \texttt{<think>} \\
\texttt{Okay, let's start by understanding the problem. I need to divide 2 chalks, 2 shoes, and 2 fruits between me and another player. My values are chalk:4, shoe:1, fruit:0. The other player's values are secret. The goal is to maximize the total score (mine plus theirs) by making complementary proposals. If the proposals are complementary, we get our respective scores; otherwise, both get zero.} \\
\\ 
\texttt{First, I should consider what the other player might value. Since my values are high for chalks and low for fruits, maybe the other player has the opposite. If they value fruits highly, then I can take the chalks and shoes, and they take the fruits. That would maximize the total because I get 2*4=8 for chalks and maybe 1 for a shoe if needed, while they get 2*their value for fruits. But I don't know their values.} \\
\\ 
\texttt{I need to communicate with them to find out their preferences. Let's start by asking them what they value. But the messages are direct, so I can't include any meta{-}text. So my first message should be something like, "What do you value most? Let me know so we can split fairly." } \\
\\ 
\texttt{If they respond that they value fruits, then I can take the chalks and shoes. If they value shoes, maybe they want shoes and I take chalks and fruits. But since I don't value fruits, maybe they do. Alternatively, if they value chalks, we might have to split them, but that could lower the total.} \\
\\ 
\texttt{Wait, the key is to maximize the total. So if I can take the items I value most and they take theirs, that's best. But without knowing their values, I need to make an educated guess. Since I have zero for fruits, it's possible they value them. Let me propose taking all the chalks and shoes, leaving fruits for them. But I need to check if that's possible.} \\
\\ 
\texttt{Alternatively, maybe they also value chalks. If we both want chalks, we might have to split them. But if I take 2 chalks, they can take the rest. Let me think of possible scenarios.} \\
\\ 
\texttt{If I propose {[}Proposal: 2 chalk, 2 shoe{]}, that would use up 4 items, leaving 2 fruits. If the other player takes 2 fruits, that's complementary. But if they need more} \\
            }
        }
    }
     \\ \\

    \theutterance \stepcounter{utterance}  
    & & & \multicolumn{2}{p{0.3\linewidth}}{
        \cellcolor[rgb]{0.9,0.9,0.9}{
            \makecell[{{p{\linewidth}}}]{
                \texttt{\tiny{[GM$|$GM]}}
                \texttt{proposal} \\
            }
        }
    }
    & & \\ \\

    \theutterance \stepcounter{utterance}  
    & & & \multicolumn{2}{p{0.3\linewidth}}{
        \cellcolor[rgb]{0.9,0.9,0.9}{
            \makecell[{{p{\linewidth}}}]{
                \texttt{\tiny{[GM$|$GM]}}
                [2, 2, 0]
            }
        }
    }
    & & \\ \\

    \theutterance \stepcounter{utterance}  
    & & \multicolumn{4}{p{0.6\linewidth}}{
        \cellcolor[rgb]{0.9,0.9,0.9}{
            \makecell[{{p{\linewidth}}}]{
                \texttt{\tiny{[P1$\langle$GM]}}
                \texttt{The other player has submitted a secret proposal. It is now your turn to submit a proposal of your own.} \\
            }
        }
    }
    & \\ \\

    \theutterance \stepcounter{utterance}  
    & \multicolumn{4}{p{0.6\linewidth}}{
        \cellcolor[rgb]{0.9,0.9,0.9}{
            \makecell[{{p{\linewidth}}}]{
                \texttt{\tiny{[P1$\rangle$GM]}}
                \texttt{<think>} \\
\texttt{Okay, the other player has already submitted a secret proposal, so I need to come up with my own. Let me think about what they might have proposed. Since I don't know their value function, I have to make an educated guess.} \\
\\ 
\texttt{My values are chalk:4, shoe:1, fruit:0. So I want as many chalks as possible and shoes are less valuable. Fruits are worthless to me. The other player's values are unknown, but since the goal is to maximize the total score (mine plus theirs), I should aim for a division where both of us get items we value highly.} \\
\\ 
\texttt{The available items are 2 chalks, 2 shoes, 2 fruits. Let's consider possible splits. If I take both chalks, that's 8 points for me. Then, the remaining items are 2 shoes and 2 fruits. If the other player values shoes and fruits, maybe they can take those. But if they value chalks too, there's a conflict. However, since I don't know their values, I need to assume they might not value the same things I do.} \\
\\ 
\texttt{Alternatively, maybe a balanced approach. Suppose I take 1 chalk and 1 shoe. That gives me 4 + 1 = 5 points. Then the other player could take the remaining 1 chalk, 1 shoe, and 2 fruits. If their values are different, say they value shoes more, then they might get more points. But without knowing their values, it's tricky.} \\
\\ 
\texttt{Wait, the key is that our proposals need to be complementary. If I take 2 chalks and 1 shoe, that's 3 items. The other player could take 1 shoe and 2 fruits. That uses all items. But if their proposal is taking 1 shoe and 2 fruits, and mine is 2 chalks and 1 shoe, then it's complementary. Let me check the counts: 2 chalks (me) + 0 (them) = 2. Shoes: 1 (me) + 1 (them) = 2. Fruits: 0 (me) + 2 (them) = 2. Yes, that works.} \\
\\ 
\texttt{So my proposal would be {[}Proposal: 2 chalk, 1 shoe{]}. That way, I maximize my points (2*4 + 1*1 = 9) and leave the remaining items for the other player} \\
            }
        }
    }
    & & \\ \\

    \theutterance \stepcounter{utterance}  
    & & & \multicolumn{2}{p{0.3\linewidth}}{
        \cellcolor[rgb]{0.9,0.9,0.9}{
            \makecell[{{p{\linewidth}}}]{
                \texttt{\tiny{[GM$|$GM]}}
                \texttt{proposal} \\
            }
        }
    }
    & & \\ \\

    \theutterance \stepcounter{utterance}  
    & & & \multicolumn{2}{p{0.3\linewidth}}{
        \cellcolor[rgb]{0.9,0.9,0.9}{
            \makecell[{{p{\linewidth}}}]{
                \texttt{\tiny{[GM$|$GM]}}
                [2, 1, 0]
            }
        }
    }
    & & \\ \\

    \theutterance \stepcounter{utterance}  
    & & & \multicolumn{2}{p{0.3\linewidth}}{
        \cellcolor[rgb]{0.9,0.9,0.9}{
            \makecell[{{p{\linewidth}}}]{
                \texttt{\tiny{[GM$|$GM]}}
                [[2, 1, 0], [2, 2, 0]]
            }
        }
    }
    & & \\ \\

\end{supertabular}
}

\end{document}
